\documentclass[12pt]{article}
\usepackage{graphicx}
\usepackage{mathtools,amsfonts,amsthm}
\usepackage{lipsum}
\usepackage{enumitem}% http://ctan.org/pkg/enumitem

\begin{document}



\newpage
\section{INTRODUCTION}
Factor Analysis is a statistical technique used to uncover the underlying sections or dimensions within a set of observed variables. This technique extracts maximum common variance from all variables and puts them into a common score. The allure of Factor analysis over other approaches to the study of hypothetical constructs is its capacity of testing detailed hypothesis in a inferential mode. The major concern of Factor analysis is modelling factors, sometimes referred to as latent variables. It determines the extent to which each variable in the dataset is associated with a common factor. Factors are influences which are not directly measured, but account for commonality among set of measurements.

\subsection{background Information}
Factor analysis is a powerful data reduction technique that enables researchers to investigate concepts that cannot easily be measured directly. They are able to test the hypothesis that a relationship between observed variables and their underlying constructs exists.It helps researchers and clinicians gain insights into the underlying constructs that contribute to mental health conditions and is useful when dealing with complex psychological phenomena that cannot be directly observed.Researchers typically start with a pool of variables (eg.symptoms and self report scales) thought to be relevant to a specific mental health construct.FA helps identify the subset of indicators that best present the underlying factors and eliminate redundant variables.Factor extraction involves identifying the underlying factors by examining the pattern of correlations among the observed variables which can be done using Principal Component Analysis (PCA).Once the factors are extracted, researchers employ factor rotation techniques (varimax, oblique rotation) to simplify the interpretation of the factor structure.

\subsection{Statement of the problem}
\noindent Many reseachers in psychology and social science are forced with the problem of comparing latent constucts that are not directly observable.\\

\noindent The latent constructs are usually measured using questionnaires, having different scales that reflect different underlying latent variables.The differences to these underlying constructs are tested via scale means.\\

\noindent Factory Analysis has become the standard to identify and interpret these underlying constructs.


\subsection{Objective of the study}
\subsubsection*{General Objectives}
To demonstrate the use of FA in reducing the dimensionality of a large set of observed variables by identifying the underlying factors that explains the patterns of correlations among variables.

\subsubsection*{Specific Objective}
\begin{enumerate}
\item To identify the latent factors contributing to presence and depression severity
\item To assess the reliability of the Factor model in explaining depression severity.
\item To show how strongly every factor affects depression severity.
\end{enumerate}

\subsection{Methodology}
This study employed factor analysis to explore the underlying factor structure in a dataset. The sample size includes 247 participants between the age of 18- 35 who took part in an online depression test, and completed a 9-item major depression screener.

\subsection{Assumptions}
\noindent Adequate sample size: It is assumed that the sample size is sufficient to provide reliable estimates and accurate parameter estimates. In general, larger sample sizes tend to yield more stable and robust factor analysis results.\\

\noindent Linearity: Factor analysis assumes that the relationships between variables are linear. Nonlinear relationships may lead to biased factor loadings and inaccurate factor structures.\\

\noindent Independence of observations: The assumption of independence assumes that each observation in the dataset is unrelated to other observations. Violations of this assumption, such as dependence or clustering, may affect the factor analysis results.\\

\noindent Homogeneity of variance: The assumption of homogeneity of variance assumes that the variables have similar variances across the range of factor scores. Violations of this assumption may lead to distorted factor loadings and biased factor solutions.\\

\subsection{justfication of the Study}
Existing gaps in the field of research and psychology in the mental health sector to determine relationships between latent variables makes the basis of this research paper.Researchers lacking the knowledge of FA are faced with the problem of latent modelling.\\
In this paper we discuss one of the tools of measuring latent variables showing the importance of this research to the researchers dealing with latent variable modelling. 

\end{document}
